\section{Первая программа}
\subsection{Исходный код}
Исходный код первой программы представлен в листинге~\ref{lst:prog-1}.

\begin{listing}[Исходный код первой программы]{lst:prog-1}{C++}
#include <fcntl.h>
#include <stdio.h>

int main()
{
    int fd = open("alphabet.txt", O_RDONLY);

    FILE *fs1 = fdopen(fd, "r");
    char buff1[20];
    setvbuf(fs1, buff1, _IOFBF, 20);

    FILE *fs2 = fdopen(fd, "r");
    char buff2[20];
    setvbuf(fs2, buff2, _IOFBF, 20);

    int flag1 = 1, flag2 = 2;
    while (flag1 == 1 || flag2 == 1)
    {
        char c;
        flag1 = fscanf(fs1, "%c", &c);
        if (flag1 == 1)
        {
            fprintf(stdout, "%c", c);
        }
        flag2 = fscanf(fs2, "%c", &c);
        if (flag2 == 1)
        {
            fprintf(stdout, "%c", c);
        }
    }
    return 0;
}
\end{listing}

\subsection{Результат работы}
Результат работы первой программы представлен на рисунке~\ref{img:1.png}.
\jimg{1.png}{Результат работы первой программы}

\subsection{Анализ полученного результата}
% Системный вызов open открывает файл \code{alphabet.txt} в режиме <<только для чтения>> и 
% создаёт структуру \code{file}. Ссылка на данную структуру помещается в таблицу файловых 
% дескрипторов процесса. Поскольку дескрипторы 0, 1 и 2 заняты стандартными потоками 
% ввода-вывода, системный вызов возвращает наименьший свободный дескриптор --- 3, 
% который присваивается переменной \code{fd}.

% Библиотечной функции \code{fdopen} передается дескриптор \code{fd} и режим доступа 
% <<только для чтения>>. Функция возвращает указатель на инициализированную 
% структуру FILE, где полю \_fileno присваивается значение fd.
% Функция setvbuf задаёт размер буфера 20 байт для fs1 и fs2.

% Функция fscanf для fs1 сначала считывает 20 символов (<<a>>---<<t>>) из 
% файла "alphabet.txt" в буфер, при этом значение поля f\_pos структуры file 
% увеличивается на 20. Затем из буфера извлекается первый символ (<<a>>) и 
% записывается в переменную c, а указатель \_IO\_read\_ptr в структуре fs1 
% сдвигается на 1 байт. Вызов fprintf выводит этот символ в стандартный поток вывода.

% Так как fs1 и fs2 ссылаются на один и тот же открытый файл, где значение 
% f\_pos равняется 20, то при вызове fscanf для fs2 в буфер считываются 
% оставшиеся 6 символов (<<u>>---<<z>>) из файла "alphabet.txt". 
% Первый символ из буфера fs2 (<<u>>) записывается в переменную c и 
% выводится в стандартный поток вывода.

% В последующих итерациях цикла поочерёдно будут выводиться символы из буферов fs1 и fs2.
% Когда будет достигнут конец буфера fs2 (сравняются указатели \_IO\_read\_ptr и \_IO\_read\_end), 
% то программа продолжит выводить оставшиеся символы только из буфера fs1 
% (вплоть до t), пока не будет достигнут конец буфера fs1.

В программе выполняется системный вызов open(), который открывает целевой 
файл в режиме <<только для чтения>> (O\_RDONLY). В результате этого в таблице 
файловых дескрипторов процесса резервируется наименьший свободный индекс 
(как правило, 3, поскольку дескрипторы 0, 1 и 2 заняты стандартными потоками 
ввода-вывода), а в системной таблице открытых файлов ядра создается 
единственная структура struct file.

Далее библиотечная функция fdopen() вызывается дважды, связывая полученный 
системный дескриптор со структурами FILE (указатели fs1 и fs2). С помощью 
функции setvbuf() для обоих файловых потоков в адресном пространстве 
пользователя выделяются буферы buff1 и buff2 размером 20 байт и устанавливается 
режим полной буферизации (\_IOFBF).

При первом вызове библиотечной функции fscanf() для указателя fs1 происходит 
чтение файла: локальный буфер buff1 заполняется первыми 20 
символами исходного файла (от <<a>> до <<t>>). При этом указатель текущей 
позиции f\_pos в структуре struct file открытого файла увеличивается на 20. Из 
буфера извлекается первый символ <<a>>, а внутренний указатель \_IO\_read\_ptr 
сдвигается на 1 байт.

Затем выполняется первый вызов fscanf() для указателя fs2. Поскольку оба 
файловых потока FILE опираются на один и тот же файловый дескриптор и, 
следовательно, на общую структуру struct file, ядро начинает чтение с 
текущей позиции f\_pos = 20. В буфер buff2 считываются оставшиеся в файле 
6 символов (от <<u>> до <<z>>), после чего первый символ из считанных (<<u>>) 
передается в переменную и выводится на экран.

В последующих итерациях цикла программа поочередно извлекает символы из 
пользовательских буферов buff1 и buff2. Системные вызовы при этом не 
производятся, так как данные уже находятся в оперативной памяти процесса. 
Как только буфер buff2 пустеет (указатели \_IO\_read\_ptr и \_IO\_read\_end 
сравниваются), программа продолжает последовательно выводить оставшиеся 
символы из буфера buff1, пока тот также не будет исчерпан.

\subsection{Схема связи структур}
\jimg{1.pdf}{Схема связи структур в первой программе}
